\subsection*{数据搬运：分块搬运引擎}

逐元素写法使用标量索引：选一个位置、读输入、计算、写输出。

分块实现则先选择区域，将其暂存到计算使用的存储中，再在那里复用数据。Croqtile 直接表达这次搬运，也能将属于搬运过程的变换写入同一语句：

\begin{itemize}[leftmargin=1.6em]
\item \co@dma.copy@ 原样搬运一个分块。
\item \co@dma.transp<...>@ 在搬运的同时置换维度。
\item \co@dma.pad<...>@ 扩大分块并填充边缘。
\end{itemize}

这一模型的核心语法和规则，决定了 DMA 代码能否正确、高效地转换为目标代码。

\begin{figure}[htbp]\centering
\includegraphics[width=\linewidth,height=0.56\textheight,keepaspectratio]{f04-fig1_element_vs_block_light.png}
\caption{标量索引逐个读取元素；显式分块搬运先选择分块，再将其复制到 shared 存储。}
\end{figure}
